\documentclass[aspectratio=169,11pt]{beamer}

\usetheme{Madrid}
\usecolortheme{default}
\usefonttheme{professionalfonts}
\setbeamertemplate{navigation symbols}{}
\setbeamertemplate{footline}[frame number]

\usepackage{amsmath, amssymb, booktabs, graphicx}

\newcommand{\E}{\mathbb{E}}
\newcommand{\1}{\mathbf{1}}
\newcommand{\Prb}{\mathbb{P}}

\title{Equilibrium Structural Work}
\subtitle{Search, Spatial, Industry/Market, and Policy Counterfactuals \\ Empirical Methods --- Week 8}
\author{Teaching deck draft}
\date{}

\begin{document}

\begin{frame}
  \titlepage
\end{frame}

\begin{frame}{Week 8 question}
\begin{itemize}
    \item When do individual decisions interact through markets, prices, firms, or locations?
    \item What does equilibrium structure buy that partial equilibrium cannot?
    \item Which moments discipline the equilibrium primitives?
    \item How should policy counterfactuals and welfare be interpreted once markets adjust?
\end{itemize}

\medskip
Equilibrium structure is worth the cost when incidence, allocation, or welfare depends on endogenous market response.
\end{frame}

\begin{frame}{When do we need equilibrium structure?}
\small
\begin{tabular}{p{0.28\linewidth} p{0.31\linewidth} p{0.29\linewidth}}
\toprule
Partial-equilibrium risk & Equilibrium object & Applied examples \\
\midrule
Other agents' environments change & wages, prices, rents, vacancies & minimum wages, UI, housing policy \\
Matching changes & tightness, sorting, queues & job search, firm recruiting, platforms \\
Locations adjust & migration, commuting, housing costs & place policy, transit, local shocks \\
Firms respond strategically & markups, entry, pass-through & mergers, taxes, technology shocks \\
Flows reallocate & bilateral frictions and shares & trade, migration, commuting, applications \\
\bottomrule
\end{tabular}

\[
\frac{\partial E_i}{\partial a_j} \neq 0
\quad \text{for many } i \neq j
\]
\end{frame}

\begin{frame}{Search and matching equilibrium}
\small
\[
\theta = \frac{v}{u},
\qquad
f(\theta)=\frac{m(u,v)}{u},
\qquad
q(\theta)=\frac{m(u,v)}{v}
\]

\[
\kappa = q(\theta)J
\]

\[
W-U=\beta S,
\qquad
J=(1-\beta)S
\]

\begin{itemize}
    \item Observed objects: unemployment, vacancies, transitions, separations, wages.
    \item Latent objects: matching technology, vacancy cost, surplus, bargaining or wage posting.
    \item Counterfactual gain: unemployment, vacancy, and wage incidence after firms respond.
    \item Main risk: matching and wage-setting assumptions carry the result.
\end{itemize}
\end{frame}

\begin{frame}{Spatial equilibrium}
\small
\[
V_{ij}
=
w_j-r_j-\tau_{ij}+A_j+\varepsilon_{ij},
\qquad
s_{ij}
=
\frac{\exp(V_{ij}/\sigma)}
{\sum_{\ell}\exp(V_{i\ell}/\sigma)}
\]

\[
N_j(w,r,A,\tau)=H_j(r_j;\eta_j),
\qquad
L_j(w,r,A,\tau)=D_j(w_j,Z_j)
\]

\[
dV_j = dw_j - dr_j + dA_j - d\tau_j
\]

\begin{itemize}
    \item Wages are not welfare when rents, amenities, and commuting costs adjust.
    \item Housing supply and mobility frictions govern incidence.
    \item Fit should include wages, rents, population or employment, and flows.
\end{itemize}
\end{frame}

\begin{frame}{Industry and market equilibrium}
\small
\[
q_j = D_j(p,x,\xi;\theta_d)
\]

\[
p_j - mc_j
=
-\left(
\frac{\partial D_j(p,x,\xi;\theta_d)}{\partial p_j}
\right)^{-1}
\;D_j(p,x,\xi;\theta_d)
\]

\[
\frac{p_j-mc_j}{p_j}=-\frac{1}{\varepsilon_{jj}},
\qquad
(p^\ast,q^\ast,L^\ast)=\Psi(\theta_d,\theta_c,\Omega,Z)
\]

\begin{itemize}
    \item Demand and conduct imply prices, markups, output, and labor incidence.
    \item Useful for taxes, mergers, platform pricing, market power, and pass-through.
    \item Validation requires credible substitution patterns and price-cost implications.
\end{itemize}
\end{frame}

\begin{frame}{Policy counterfactuals and welfare}
\small
\[
x^\ast(\tau)=\Phi(x^\ast(\tau),\tau;\theta)
\]

\[
W(\tau)
=
\sum_i \omega_i U_i(x^\ast(\tau),\tau;\theta)
+
\sum_f \Pi_f(x^\ast(\tau),\tau;\theta)
+
R(x^\ast(\tau),\tau)
\]

\[
\Delta W = W(\tau_1)-W(\tau_0)
\]

\begin{itemize}
    \item The fixed point collects wages, vacancies, rents, prices, shares, entry, or flows.
    \item Welfare depends on distributional weights, profits, rents, and fiscal terms.
    \item Report counterfactual distance from support and sensitivity to closure assumptions.
\end{itemize}
\end{frame}

\begin{frame}{Computational and identification burdens}
\small
\begin{tabular}{p{0.24\linewidth} p{0.34\linewidth} p{0.30\linewidth}}
\toprule
Approach & Intuition & Risk \\
\midrule
Nested fixed point & solve equilibrium inside parameter search & slow and opaque \\
MPEC & estimate with equilibrium constraints directly & numerical delicacy \\
Calibration-estimation hybrid & estimate some objects, calibrate others & headline result may rest on calibration \\
Simulation & approximate heterogeneous equilibrium response & seed, tolerance, and approximation error \\
\bottomrule
\end{tabular}

\[
\widehat m^{data}
\Longleftrightarrow
m^{model}(\theta,x^\ast(\theta))
\]

\begin{itemize}
    \item Name which moments identify which primitives.
    \item Show targeted and untargeted fit.
    \item Report solver, tolerances, starts, convergence, and sensitivity.
\end{itemize}
\end{frame}

\begin{frame}{Gravity-style equilibrium models}
\small
\[
M_{ij}
=
A_iB_j\tau_{ij}^{-\sigma}
\]

\begin{itemize}
    \item The object is a bilateral flow: trade, migration, commuting, applications, or worker allocation.
    \item The structural object is a friction or resistance term that governs reallocation.
    \item Gravity belongs here when flows close an equilibrium with wages, prices, rents, or employment.
    \item It should remain subordinate to the broader question: how does equilibrium incidence work?
\end{itemize}

\medskip
Diagnostic question: does the model only describe flows, or does it solve the market response that matters for welfare?
\end{frame}

\begin{frame}{Research Lab design}
\begin{columns}[T,onlytextwidth]
\begin{column}{0.32\linewidth}
\textbf{Reproduce}
\begin{itemize}
    \item Hsieh-Moretti-style spatial equilibrium
    \item wages, rents, population
    \item housing relaxation counterfactual
    \item output and welfare table
\end{itemize}
\end{column}
\begin{column}{0.32\linewidth}
\textbf{Diagnose}
\begin{itemize}
    \item housing elasticity
    \item mobility elasticity
    \item targeted vs closure objects
    \item validation memo
\end{itemize}
\end{column}
\begin{column}{0.32\linewidth}
\textbf{Transfer}
\begin{itemize}
    \item BLP/Nevo-style market exercise
    \item demand elasticity
    \item markup recovery
    \item cost-shock pass-through
\end{itemize}
\end{column}
\end{columns}
\end{frame}

\begin{frame}{Takeaways}
\begin{itemize}
    \item Equilibrium structure is a research choice, not a badge of sophistication.
    \item Search, spatial, and market models differ in economics but share an empirical architecture.
    \item The model must say which equilibrium margin matters and which moments discipline it.
    \item Counterfactuals are credible only when fit, support, computation, and closure are transparent.
    \item The best equilibrium papers make the gains and burdens visible in the same language.
\end{itemize}
\end{frame}

\end{document}
